%% ============================================================================
%% Section: Related Work
%% ============================================================================
\section{Related Work}\label{sec:related}

\noindent\textbf{Reachability indexing.} Answering reachability without
per-query traversal is a long line of database research, from 2-hop
labeling~\cite{cohen2003reachability} through scalable interval and label
indexes such as GRAIL~\cite{yildirim2010grail}, FERRARI~\cite{seufert2013ferrari},
TF-Label~\cite{cheng2013tflabel}, pruned landmark
labeling~\cite{yano2013pruned}, IP$^+$~\cite{wei2014ip}, and
BFL~\cite{su2017bfl}, to DAG reduction~\cite{zhou2017dagreduction} and
O'Reach~\cite{hanauer2022oreach}; a recent survey covers the
field~\cite{reachabilitysurvey2025}. These indexes answer point-to-point
reachability on graphs with millions of nodes under a label-size budget. Our
ancestor and descendant bitsets are a deliberately simple point in this
space, a full transitive closure packed into machine words, chosen for two
reasons: causal DAGs are small enough for the closure to fit, and the
operators above reachability (confounders, forbidden sets, canonical
adjustment sets, minimal separators) need whole ancestor and descendant sets,
which bitsets deliver as one word-parallel operation. We bring a classical
reachability trade-off to a workload that has not previously been treated as
an indexed DBMS workload, and we maintain the closure incrementally, which
label-based indexes typically do not.

\noindent\textbf{Property-graph query languages.} Cypher and the GQL and
SQL/PGQ standards~\cite{deutsch2022gql,francis2023gpc} express graph pattern
matching and variable-length paths declaratively; our language compiles to
Cypher plus procedures rather than replacing it. What these languages do not
provide is conditioning-aware d-separation or adjustment-set construction as
first-class operators: these are not path patterns, and expressing them in
Cypher means client-side orchestration of many traversals. Our earlier vision
paper~\cite{pachera2025whatif} argued for causal operators over property
graphs and sketched extraction, do-calculus, and view-maintenance directions
with a proof of concept on Neo4j; the present paper contributes what that
vision left open, namely a concrete language, a physical index, the
query-processing algorithms over it, incremental maintenance, and an
evaluation, and a companion position paper~\cite{qiu2026causal} places such a
store in a broader causal data-management ecosystem.

\noindent\textbf{Causality in databases.} Database research has studied
causal inference over relational data~\cite{salimi2018hypdb,salimi2020carl},
hypothetical what-if and how-to queries with a probabilistic causal
model~\cite{galhotra2022hyper}, causal data
integration~\cite{youngmann2023integration}, and causality as a tool for
explaining query answers~\cite{meliou2014causality}. In these works the
causal model is an input to query answering over the data; in ours the causal
DAG is the stored object and the operators are structural queries over it.
The two are complementary: an identification query answered by Causal4j
selects the adjustment set that a relational causal-inference engine then
uses to estimate an effect.

\noindent\textbf{Causal tooling and algorithms.} Causal analysis is dominated
by standalone libraries (DoWhy~\cite{sharma2020dowhy},
dagitty~\cite{textor2016dagitty}, pgmpy~\cite{ankan2024pgmpy},
causal-learn~\cite{zheng2024causallearn}, Ananke~\cite{lee2023ananke}) that
implement d-separation, the adjustment criterion, and identification over an
in-memory graph. None is DBMS-resident: each requires exporting the DAG,
computes in a separate process, and discards its representation, so none
offers durable shared access, in-database querying, or maintenance under a
changing graph. The algorithms we build on come from this community:
Bayes-Ball~\cite{shachter1998bayesball} and the ancestral restriction of
d-separation~\cite{geiger1990identifying}, the adjustment criterion and its
completeness~\cite{shpitser2010validity,perkovic2018complete}, and the
separator framework with its constructive identifiability test and
polynomial-delay enumeration~\cite{textor2011adjustment,vanderzander2014constructing,vanderzander2019separators,vanderzander2020minimal}.
We reuse their semantics (indeed we validate against dagitty in
Section~\ref{sec:eval}) and relocate the computation inside the store, where
the sets these algorithms consume are precomputed.

\noindent\textbf{In-DBMS analytics and incremental maintenance.} Pushing
analytics into the database through procedures and user-defined functions is
established practice; our contribution is the physical design and query
processing for causal operators, not the extension mechanism. Keeping the
index consistent under updates connects to incremental view
maintenance~\cite{gupta1995maintenance} and to dynamic transitive
closure~\cite{demetrescu2004dynamic}: our insertion patch is an instance of
maintaining a materialized closure under edge insertion, and our cone
recompute bounds the non-local case of deletion by the ancestors and
descendants of the deleted edge, with recompilation as the fallback.
